\chapter{Mutation with PIT} \label{cap:mutation}

\section{Mutation Analysis Setup}

Program mutation testing was executed with PIT over class \texttt{sut.TST}.
The comparison goal was restricted to method \texttt{longestPrefixOf}, using one criterion-specific class per run:
Edge-Pair, Prime Path, and Logic-based (CDC).
Each run selected the relevant test class through \texttt{-DtargetTests=...}, so mutation outcomes could be compared under controlled test-selection conditions.
In this project configuration, PIT summaries may still include mutants from multiple methods in \texttt{sut.TST};
therefore, comparison values were computed by filtering XML entries with \texttt{mutatedMethod = longestPrefixOf}.

\section{Execution Commands}

The following command sequence can be used to reproduce the PIT runs used in this report.
Because reports are generated in \texttt{target/pit-reports}, copying \texttt{mutations.xml} after each run keeps per-criterion artifacts separate.

\begin{verbatim}
# Edge-Pair suite
mvn -DtargetTests=sut.coverage.longestprefix.edgepair.LongestPrefixOfEdgePairTest \
    org.pitest:pitest-maven:mutationCoverage
cp target/pit-reports/mutations.xml target/pit-reports/mutations-edgepair.xml

# Prime Path suite
mvn -DtargetTests=sut.coverage.longestprefix.primepath.LongestPrefixOfPrimePathTest \
    org.pitest:pitest-maven:mutationCoverage
cp target/pit-reports/mutations.xml target/pit-reports/mutations-primepath.xml

# Logic-based (CDC) suite
mvn -DtargetTests=sut.coverage.longestprefix.logic.LongestPrefixOfConditionDecisionTest \
    org.pitest:pitest-maven:mutationCoverage
cp target/pit-reports/mutations.xml target/pit-reports/mutations-cdc.xml
\end{verbatim}

\section{Comparison for \texttt{longestPrefixOf}}

The project report \texttt{longestPrefixOf-mutation-comparison.md} was used to consolidate results per criterion.
Results below were refreshed from PIT reruns on \texttt{2026-04-19}, using
\texttt{target/pit-reports/mutations-edgepair.xml},
\texttt{target/pit-reports/mutations-primepath.xml}, and
\texttt{target/pit-reports/mutations-cdc.xml}.

\begin{table}[H]
	\centering
	\begin{tabular}{|l|r|r|r|r|r|}
		\hline
		Criterion & Total Mutants & Killed & Survived & No Coverage & Mutation Score \\
		\hline
		Edge-Pair & 13 & 13 & 0 & 0 & 100\% \\
		Prime Path & 13 & 13 & 0 & 0 & 100\% \\
		Logic-based (CDC) & 13 & 13 & 0 & 0 & 100\% \\
		\hline
	\end{tabular}
	\caption{Mutation comparison for \texttt{TST.longestPrefixOf}}
\end{table}

\begin{figure}[H]
	\centering
	\includegraphics[width=0.95\textwidth]{mutation/images/mutationresult.png}
	\caption{PIT mutation report evidence used in the criterion comparison}
	\label{fig:pit-mutation-evidence}
\end{figure}

Although final mutation scores are identical, PIT's \texttt{killingTest} field indicates that each criterion kills mutants through different test obligations.

\begin{table}[H]
	\centering
	\scriptsize
	\setlength{\tabcolsep}{3pt}
	\renewcommand{\arraystretch}{1.15}
	\begin{tabular}{|p{2.1cm}|p{4.9cm}|p{6.6cm}|}
		\hline
		Criterion & Dominant killing tests (from PIT XML) & Interpretation \\
		\hline
		Edge-Pair & \texttt{traversalTakesLeftAnd}\newline \texttt{RightEdges} (6);\newline \texttt{equalEdgeWithoutValue}\newline \texttt{DoesNotUpdateLength} (3) & Most kills come from edge-transition combinations and value-update edge cases. \\
		\hline
		Prime Path & \texttt{primePath\_rightThen}\newline \texttt{EqualThenMid} (7);\newline \texttt{primePath\_exitsWhen}\newline \texttt{SearchFallsOffTrie} (4) & Most kills come from full-path traversal obligations and path-exit behavior. \\
		\hline
		Logic-based (CDC) & \texttt{cdc\_whileCondition\_true}\newline \texttt{BecauseBothClausesTrue} (8);\newline \texttt{cdc\_whileCondition\_false}\newline \texttt{BecauseLeftClauseFalse} (3) & Most kills come from independent truth/falsity control of \texttt{x != null} and \texttt{i < query.length()}. \\
		\hline
	\end{tabular}
	\normalsize
	\caption{Dominant mutant-killing tests by criterion}
\end{table}

\section{Improvements Performed}

Historical pre-improvement PIT XML was not archived in this workspace, so a fully reproducible baseline table with exact earlier counts cannot be reconstructed.
What is reproducible in the current workspace is the improvement action and the final rerun state.
Two targeted guard tests were added to close entry-guard gaps in \texttt{longestPrefixOf}:
\texttt{primePath\_entryGuardEmptyQuery} in \texttt{LongestPrefixOfPrimePathTest} and
\texttt{cdc\_entryGuardEmptyQuery} in \texttt{LongestPrefixOfConditionDecisionTest}.
Current reruns confirm that all three criteria now kill all 13 \texttt{longestPrefixOf} mutants.

\section{Discussion}

The final result shows that all three criteria killed the same set of \texttt{longestPrefixOf} mutants.
Even with identical final scores, the criteria provide complementary evidence: edge-pair tests emphasize local transition correctness, prime-path tests stress deeper path behavior, and CDC targets decision sensitivity.
Using these perspectives together increases confidence in both control-flow and logic-level correctness.

For completeness, the remaining non-\texttt{longestPrefixOf} surviving mutants are boundary changes in
\texttt{collect} (lines 219 and 229), used by \texttt{keysThatMatch}.
They are likely equivalent, since changing \texttt{<} to \texttt{<=} and \texttt{>} to \texttt{>=} at those points only explores an extra null branch in wildcard traversal and does not change the returned set of matched keys.